% from ega3-1-fr.tex
\oldpage[0\textsubscript{III}]{52}
En outre, on a des diagrammes commutatifs
\[
\label{0.12.1.4.2-fr}
\tag{12.1.4.2}
\xymatrix{
\check H^p(\mathfrak U,\mathcal F)\ar[r]^{\theta_p}\ar[d]
&\check H^p(\mathfrak U',f^*(\mathcal F))\ar[d]\\
H^p(X,\mathcal F)\ar[r]_{\theta_p}
&H^p(X',f^*(\mathcal F))
}
\]
où les flèches verticales sont les homomorphismes canoniques de (G, II, 5.2).
Pour établir la commutativité de
(\hyperref[0.12.1.4.2-fr]{12.1.4.2}), considérons le complexe de faisceaux de
cochaînes (alternées) de $\mathcal F$ relatives à $\mathfrak U$,
$\mathcal C^\bullet(\mathfrak U,\mathcal F)$, tel que
$\Gamma(X,\mathcal C^\bullet(\mathfrak U,\mathcal F))=
\check C^\bullet(\mathfrak U,\mathcal F)$ (G, II, 5.2). Les homomorphismes
canoniques $\Gamma(V,\mathcal F)\to
\Gamma(f^{-1}(V),\psi^*(\mathcal F))$ définissent alors un $\psi$-morphisme
$\mathcal C^\bullet(\mathfrak U,\mathcal F)\to
\mathcal C^\bullet(\mathfrak U',\psi^*(\mathcal F))$, et on a, avec les
notations ci-dessus, un diagramme commutatif
\[
\xymatrix{
\Gamma(X,\mathcal C^\bullet(\mathfrak U,\mathcal F))\ar[r]\ar[d]
&\Gamma(X',\mathcal C^\bullet(\mathfrak U',\psi^*(\mathcal F)))\ar[d]\\
\Gamma(X,\mathcal L^\bullet)\ar[r]
&\Gamma(X',\psi^*(\mathcal L^\bullet))\ar[r]
&\Gamma(X',\mathcal L'^\bullet)
}
\]
qui, en passant à la cohomologie, donne des diagrammes commutatifs
\[
\xymatrix{
\check H^p(\mathfrak U,\mathcal F)\ar[r]\ar[d]
&\check H^p(\mathfrak U',\psi^*(\mathcal F))\ar[d]\\
H^p(X,\mathcal F)\ar[r]
&H^p(X',\psi^*(\mathcal F))
}
\]
où les flèches verticales sont les homomorphismes canoniques de (G, II, 5.2).
Il suffit alors de combiner ces diagrammes avec les diagrammes commutatifs
\[
\xymatrix{
\check H^p(\mathfrak U',\psi^*(\mathcal F))\ar[r]\ar[d]
&\check H^p(\mathfrak U',f^*(\mathcal F))\ar[d]\\
H^p(X',\psi^*(\mathcal F))\ar[r]
&H^p(X',f^*(\mathcal F))
}
\]
